%% kapa: exam=1 round=29 year=2018 period=2 form=A subject=appraisal_law subject_ko=감정평가관계법규 seq=20 number=20
%% source: Q-Net 감정평가사 기출문제 (cst00309) artlSeq=5202808 file=extracted/1008/5202808/2018년 제29회 감정평가사 1차시험 기출문제/2018년 제29회 감정평가사 1차시험 기출문제/2교시 A형.pdf
%% conversion: kordoc 4.24.51 -> tools/gen1.py (auto); review=proofread-sonnet; flags=none
\begin{question}{20}
국유재산법령상 국유재산에 관한 설명으로 옳지 않은 것은?
\begin{choices}
  \item 대통령 관저와 국무총리 공관은 공용재산이다.
  \item 행정재산 외의 모든 국유재산은 일반재산이다.
  \item 행정재산의 사용 여부는「국가재정법」제6조에 따른 중앙관서의 장의 의견을 들어 기획재정부장관이 결정한다.
  \item 국가가 보존할 필요가 있다고 국토교통부장관이 결정한 재산은 보존용재산이다.
  \item 정부기업이 비상근무에 종사하는 직원에게 제공하는 해당 근무지의 구내 또는 이와 인접한 장소에 설치된 주거용 시설은 기업용재산이다.
\end{choices}
\end{question}
\answer{4}
